% !TEX root = ../main.tex
\section{Hạn chế và Thách thức}

Mặc dù kiến trúc D3QN kết hợp với Prioritized Experience Replay (PER) đã chứng minh được hiệu quả vượt trội trên môi trường Pong, hệ thống vẫn còn tồn tại một số hạn chế và thách thức kỹ thuật cần được khắc phục trong tương lai:

\begin{enumerate}
    \item \textbf{Độ phức tạp tính toán của SumTree khi Buffer lớn:} Việc sử dụng cấu trúc SumTree cho PER giúp tối ưu hóa thời gian bốc mẫu xuống $O(\log N)$. Tuy nhiên, khi áp dụng cho các môi trường phức tạp đòi hỏi dung lượng Replay Buffer khổng lồ (ví dụ $10^6$ hoặc $10^7$ transitions), việc duy trì và cập nhật cây nhị phân liên tục ở mỗi bước học vẫn tạo ra một nút thắt cổ chai (bottleneck) về hiệu suất CPU, làm chậm đáng kể quá trình phân phối dữ liệu (data pipeline) lên GPU so với việc bốc mẫu ngẫu nhiên.
    
    \item \textbf{Sự nhạy cảm với các siêu tham số (Hyperparameter Sensitivity):} Mô hình Học tăng cường sâu cực kỳ nhạy cảm với cấu hình hệ thống. Tốc độ học (Learning rate), tần suất cập nhật Target Network, và đặc biệt là hệ số bù trừ $\beta$ của thuật toán PER có ảnh hưởng sống còn đến độ ổn định của Loss. Việc tinh chỉnh (tuning) các tham số này chủ yếu dựa trên kinh nghiệm thực nghiệm (heuristics) thay vì lý thuyết toán học vững chắc, gây tốn kém tài nguyên tính toán.
    
    \item \textbf{Hiện tượng Deadly Triad:} Bất chấp các cải tiến lớn của Double DQN và Dueling Network nhằm chống lại Overestimation, hệ thống tổng thể vẫn là sự kết hợp của 3 yếu tố nguy hiểm (Deadly Triad) trong Reinforcement Learning: Function Approximation (Dùng mạng CNN để xấp xỉ), Bootstrapping (Cập nhật Q-value dựa trên ước tính của trạng thái tiếp theo), và Off-policy Learning (Sử dụng kinh nghiệm cũ từ Replay Buffer). Sự kết hợp này luôn tiềm ẩn rủi ro phân kỳ toán học (divergence) ở các game Atari phức tạp hơn nếu Loss không được kiểm soát tốt.
\end{enumerate}
